\scanpage{151}
% archive foliation: 150.
Considérons maintenant des courbes elliptiques
\inserted{$E$} avec ``\uline{rigidi}fication \unsure{différentielles|différentielle}'', définie
par une base $\varphi$ de $t_E$, permettant donc
d'identifier $t_E$ à $\mathbb C$. La donnée d'une
courbe elliptique avec rigidification différentielle
s'identifie donc à la donnée d'un \uline{réseau}
$H\subset\mathbb C$ dans $\mathbb C$. La donnée d'une courbe
elliptique avec rigidification double Torelli
stricte $+$ différentielle s'identifie donc à
la donnée de
\[
z_1,z_2\in\mathbb C^*
\quad\text{avec}\quad
\deleted{\unsure{z_2|z}}\ \tau=z_2/z_1\in\mathcal D^+
\tag{17}
\]
i.e. \deleted{la $\mathbb R$-base de $\mathbb C$}
\inserted{$z_1,z_2$ \unsure{formant|forment} une}
$\mathbb R$-base de $\mathbb C$,
qui est \uline{directe} pour l'orientation canonique
de $\mathbb C$. On fait opérer $\mathrm{Sl}(2,\mathbb Z)$ à
gauche sur la variété analytique
complexe
\[
P\mathcal D^+=\{(z_1,z_2)\mid z_1,z_2\in\mathbb C^*,\ z_2/z_1\in\mathcal D^+\}
\tag{18}
\]
par
\[
g(z_1,z_2)=(dz_1+cz_2,bz_1+az_2)
\tag{\unsure{19|18}}
\]
\margin{c'est donc, via le passage $g\mapsto\tau_\infty(g^{-1})$,
l'opération naturelle qui découle du pb modulaire}
La famille modulaire bi-rigidifiée $\unsure{\mathcal X|\mathfrak X}_{P\mathcal D^+}$
peut s'écrire
